\section{怎样写信}

我们现在再来谈谈写信时应该注意哪些问题。

\textbf{1. 掌握写信的惯用格式。}

写信与写日记或写一般文章不同，主要由于写作的目的和读者不同。它不像写日记那样，记录自己的生活经历和感受，以激励自己奋进和供日后查考，读者就是自己；它也不像一般文章那样，反映社会生活，表达自己见解，以教育人民、团结人民，打击敌人、消灭敌人，推动社会的前进，读者包括社会上的很多人。写信则是因为不能与对方口头交谈而只好进行笔谈的一种表达方式，目的在于把自己的工作、学习、生活、思想等情况告诉对方，或是要了解对方的一些情况，或是有求于对方，或是与对方进行某种联系，读者只是接信的对方。为了能使对方收到信，看懂信，看好信，首先就得使信符合惯用的格式。

书信一般由信封和信的内容两部分组成。信的内容一般要包括“称呼、问候、正文、结尾、署名、日期”六个方面，有时还可以追加“附言”。

信封左上方写收信人的地址，地址要写准、写全、写清楚，否则就可能使收信人收不到信。信封中间部分写收信人的姓名，收信人名字后面可以加上“同志收”、“先生收”，字要写得稍大一点，笔画一定要清楚、工整。信封右下侧写寄信人的地址和姓名，也要写准、写全、写清楚，以便对方回信，或者万一投递不到可以退回。在信封右上角贴上邮票。这是说的横写信封，如果是直写信封，则从右侧直写收信人地址，中间直写收信人姓名，然后在左下侧直写寄信人地址和姓名就行了。

信的内容应按写信的惯用格式来写：

第一行顶格写对收信人的称呼，称呼后面加个冒号，表示下面有话要说。一般说来，平时口头怎样称呼，信上就怎样称呼。对同学、同事、战友一般写名字，或在名字后面加“同学”、“同志”；对不熟悉的人，不仅写名字，还要写姓，姓名之后加“同志”，对长辈一般不用写名字，只写尊敬的称呼，如“爸爸”、“姑姑”；对老师、师傅或有声望的人，为了表示尊重，可以用职务名称来称呼，如“张老师”、“王师傅”、“孙大夫”、“赵主任”等。总之，称呼既要合乎彼此的关系，又要有礼貌，千万不要仅仅直呼收信人的姓名。著名的语言学家王力先生说过：“有的同志在信封上干脆写王力收，那更不好。我回信说：‘你在信内称我做尊敬的王力教授，太客气了；你在信封上写王力收，又太不客气了。’这是礼貌问题。”（《谈谈写信》）王力先生的话，很值得我们注意。

第二行开头空两格，写对收信人问候的话，可以独立成一小段。因为人们见面总要说问候的话，写信也应这样，显得很有礼貌。问候语一般用“您好！”或者“您们好！”也可用“新年好！”“你们近来一切皆好吧，我很挂念。”“你近来身体很好吧，精神一定很愉快！”等等来表示。

第三行起，是信的正文，即信的主要部分，每段起行都要空两格。写信人要说的话，要办的事，都写在这里。如果写的是回信，最好先写明“×月×日来信收悉”，接着回答来信中提出的问题或要求办的事，有什么结果告诉对方，然后再写自己想说的话。

正文完之后，用表示祝愿或敬意的话来结尾，收束全文。如“祝你幸福”，“此致敬礼”，“祝一路顺风”，“谨颂教安（编安）”等等。无论怎样写，总要使对方有一种亲切感，是由衷的祝愿。结尾这句话一般分两行写，以“祝你幸福为例”，“祝你”写在正文之后第一行空两格处，“幸福”写在下一行的顶格处，后面不必加感叹号。

署名写在结尾的话再下一行的后半行。署名的写法，和开头的称呼要对应。一般说来，对熟悉的人可以不写姓，只写名字就可以了；对不熟悉的人或对初次写信的人则要写上姓和名。署名的字体略小一点，表示礼貌。

署名的下一行写发信的年月日。

如果写完了信，又想到一些什么事情要告诉对方，可以在最后补写几句“附言”——先写一个“附”字，后面加上冒号；写完补充的话，再加上“又及”两字来说明。

\textbf{2. 正文要做到表达意思清楚明白，简洁扼要。}

王力先生说：“写信没有什么秘诀，顺着自然就是了，写信就是谈话。由于对话人相隔太远，没法子面谈，所以改为笔谈。如果我们写信仍照日常说话一样，不装模作样，不改变现代汉语的语法和词汇，就不会出毛病。”是呀，写信就是谈话，只不过是笔谈罢了。面谈时话没听明白，可以马上问；信写不明白，对方只好写信来问，这多麻烦！如果由于表达不清楚，而被对方误解，甚至误了事，后果可严重了。所以，写信要告诉对方什么事情，要求对方办什么事情，或者回答什么问题，必须说得一清二楚。这是写信至关紧要的一条。

例如，一九八五年全国高考语文试卷的作文题，要求考生以“澄溪中学学生会”的名义，给《光明日报》编辑部写一封信，反映学校附近一家工厂严重地污染环境、影响师生健康的情况，并且针对这家工厂借口——没钱、二没技术力量、三没时间解决这个问题，进行反驳，提出建议，发出呼吁。结果有的考生突如其来地指责工厂领导人无理，或者直接请求《光明日报》主持正义，或者把这家工厂的借口放在一边，而大谈环境污染的危害，或者把事件的经过说得颠三倒四。可以设想，《光明日报》编辑部接到这样的信，连事情的来龙去脉、彼此争辩的焦点是什么都搞不清楚，怎能贸然处理这件事呢？有的考生则在信中堆砌词藻，往往弄得用词不当，词不达意；有的考生则写得半文半白，语意晦涩。试想，意思都没表达清楚，你写这信还有什么作用呢？我们该记住王力先生的话，一定要用合乎规范的文风朴实的语言，把要说的意思有条不紊、清楚明白地告诉对方，千万不可华而不实，故弄玄虚，弄巧成拙。

在表意清楚的同时，还要写得简洁扼要。因为大家的学习、工作都很快，写信花时间，看信也要花时间，如果语言啰啰嗦嗦，语意拐弯抹角，说些不相干的废话，或者说些鸡毛蒜皮的话，不是使写信人浪费时间，收信人也浪费时间，影响了学习和工作么？何况说些用不着的话，还会妨碍对方对信里重要内容的理解，这也会误事的。因此，王力先生说：“写信要开门见山，不相干的话少说。”这应该是我们写信时的一条原则。像刚才提及的那道高考作文题，是要考生反映情况，表明看法，发出呼吁，偏偏有些考生一开头就说了一大堆不切题目要求的废话，把《光明日报》恭维一大通，说编辑部同志多么辛苦，贡献多么大，以前做得多么好，甚至提到登了某篇读者来信解决了问题，等等，这才说到“我们也给贵报反映一个情况，请予支持。”相反，对于相邻工厂的三个借口只说一句“统统都是站不住脚的”轻轻带过，这真是拣了芝麻、丢了西瓜了。

\textbf{3. 表达要恰当运用记叙、议论、说明、抒情等方式，并且要考虑对象，讲究礼貌。}

书信是一种带有综合性的文体。即使在叙事为主的信中，也得适当运用议论、说明或抒情；或者在议论为主的信中，也得适当运用记叙、说明或抒情。运用哪种表达方式为主，得由信的内容的性质来决定。例如代表学生会给《光明日报》编辑部写信，反映某工厂污染环境、学校提出意见又遭某工厂领导以三条“理由”拒绝的情况，这一段文字当然应以记叙为主，兼用说明，把事情经过情形简明扼要地介绍，而分析、反驳某工厂领导人所说的三条“理由”是这封信的重点部分，必然要用议论的方式进行表达，分析问题实质，指出危害，表明见解，最后发出呼吁，文字宜短，但感情却要强烈。如果一封信能恰当地运用各种表达方式，就能确切地反映事实，说明客观情况，表明自己的看法、要求或建议，也就能使收信人了解情况，决定怎么办。不然，写信的作用将会受到影响。

在写信时，要考虑对方的身份，同自己的关系以及文化程度，等等，采取相应的话语，感情真挚，实事求是地进行表达，使对方感到既尊重他，又不是假客套。这儿的关键是要掌握好分寸。礼貌要周到，不要失礼；话语要诚恳，不要做作；做到这两点，收信人就会感到愉快，信的内容也就容易被接受了。
\newpage